%% !TEX root = ../manual_gaolv5.tex
% Appendix A.
\chapter{Documentation and references}\label{chap:documentation}

\section{Documentation on GAOL}

The primary reference is this manual. The directory \file{doc} of the sources
holds the pages that give the details of GAOL v5, with the measures behind
its choices, whose content the chapters of this manual take up:

\begin{itemize}
\item \file{doc/building.md}: the CMake, autotools and meson builds, and
  their options (Chapter~\ref{chap:installation});
\item \file{doc/using.md}: the flags of interval arithmetic, GAOL from CMake
  and from pkg-config, the namespaces, the rounding direction and the
  floating-point environment (Chapters~\ref{chap:overview}
  and~\ref{chap:rounding});
\item \file{doc/three-builds.md}: what the three builds agree on, what the
  options are worth, and the compilers and options they refuse
  (Chapter~\ref{chap:installation});
\item \file{doc/tests.md}: what the programs of \file{tests} check
  (Chapter~\ref{chap:tests});
\item \file{doc/accuracy.md}: the tightness of each operation, with its
  algorithm and what the tests found (Chapter~\ref{chap:accuracy});
\item \file{doc/differences.md}: each change GAOL v5 brings to GAOL 4.2.2,
  and where it comes from (Chapter~\ref{chap:changes});
\item \file{doc/continuous-integration.md}: the systems, processors and
  compilers GAOL is built and tested on (Chapter~\ref{chap:tests});
\item \file{doc/simd.md}: what GAOL computes with vector instructions and
  with the fused multiply-add, where the time of a program goes, and the
  rounding direction by processor (Chapter~\ref{chap:computing});
\item \file{doc/math-libraries.md}: the mathematical libraries libm,
  mathlib, CRlibm and CORE-MATH, and what separates them
  (Chapter~\ref{chap:concepts});
\item \file{doc/tools.md}: the scripts of the repository and what each is
  for (Chapter~\ref{chap:tests});
\item \file{doc/compare/README.md}: GAOL compared with libieeep1788,
  filib++, PROFIL/BIAS, Solaris Studio and Boost.Interval, their special
  cases against IEEE 1788-2015, the enclosure of the elementary functions of
  GAOL and Boost.Interval, their performance, and the test suites ITF1788 and
  of Tang et al. (Chapter~\ref{chap:computing});
\item \file{examples/examples.md} and \file{examples/tutorial.md}: the
  example programs, a review of GAOL v5 and a comparison with other interval
  libraries, and a short tutorial (Chapter~\ref{chap:tests});
\item \file{3rd/README.md}: the version of CORE-MATH GAOL takes, what differs
  from it, and the fixes to propose upstream.
\end{itemize}

\noindent
The manual of GAOL~4, by Fr\'ed\'eric Goualard, is kept in \file{manual/v4}; its
PDF, \file{manual/v4/Manual\_GAOLv4.pdf}, was built once, and no build makes it
again.

\section{References}

The following articles, books and standards have inspired in some way or
another the devising of the GAOL library and the writing of this manual.

\begin{itemize}
\item \emph{IEEE Standard for Interval Arithmetic}. IEEE Std
  1788-2015~\cite{IEEE1788:2015}.
\item \emph{The CORE-MATH Project}. Alexei Sibidanov, Paul Zimmermann and
  St\'ephane Glondu~\cite{CoreMath:2022}, and
  \url{https://core-math.gitlabpages.inria.fr/}.
\item \emph{How do you compute the midpoint of an interval?} Fr\'ed\'eric
  Goualard~\cite{Goualard:TOMS2014}.
\item \emph{Fast and Correct SIMD Algorithms for Interval Arithmetic}.
  Fr\'ed\'eric Goualard~\cite{Goualard:PARA08}.
\item \emph{Interval Extensions of Multivalued Inverse Functions}.
  Fr\'ed\'eric Goualard~\cite{Goualard:2008}.
\item \emph{Interval Arithmetic Specification}. Dmitri Chiriaev and G.\:
  William Walster. Draft revised May 1998.
\item \emph{The Extended Real Interval System}. G.\: William Walster. April
  1998.
\item \emph{C++ Interval Arithmetic Programming Reference}. Sun Microsystems,
  Inc. October 2000, revision A.
\item \emph{Interval Arithmetic: From Principles to Implementation}. T.\:
  Hickey, Q.\: Ju, and M.\: H.\: van Emden~\cite{Hickey-Ju-vanEmden:2001}.
\end{itemize}

\section{Reporting bugs}

Bugs and suggestions for improvement of GAOL v5 are submitted as issues
of its repository:

\begin{center}
  \url{https://github.com/Jordan08/GAOL/issues}
\end{center}

\noindent
A report is most useful with the program that shows the problem, the
interval it gives and the one expected, the system, the processor, the
compiler and the options GAOL was built with.
